\documentclass[aspectratio=169,10pt]{beamer}

% ─── Thème et couleurs ───────────────────────────────────────────────────────
\usetheme{Madrid}
\definecolor{rtypeblue}{RGB}{0,80,150}
\definecolor{rtypedark}{RGB}{20,20,40}
\definecolor{rtypegreen}{RGB}{40,140,60}
\definecolor{rtypeorange}{RGB}{220,120,20}
\definecolor{rtypegray}{RGB}{100,100,100}
\definecolor{rtypelight}{RGB}{230,240,250}
\usecolortheme[named=rtypeblue]{structure}
\setbeamercolor{title}{fg=white,bg=rtypeblue}
\setbeamercolor{frametitle}{fg=white,bg=rtypeblue}
\setbeamercolor{block title}{bg=rtypeblue,fg=white}

% ─── Packages ────────────────────────────────────────────────────────────────
\usepackage[utf8]{inputenc}
\usepackage[T1]{fontenc}
\usepackage[french]{babel}
\usepackage{listings}
\usepackage{booktabs}
\usepackage{graphicx}
\usepackage{hyperref}
\usepackage{array}
\usepackage{adjustbox}
\usepackage{colortbl}
\arrayrulecolor{gray!50}
\renewcommand{\arraystretch}{1.3}
\usepackage{tikz}
\usetikzlibrary{positioning,arrows.meta,shapes.geometric,fit,calc,backgrounds}
\usepackage{forest}

% ─── Listings config ─────────────────────────────────────────────────────────
\lstset{
  language=C++,
  basicstyle=\ttfamily\tiny,
  keywordstyle=\color{rtypeblue}\bfseries,
  commentstyle=\color{gray},
  stringstyle=\color{red!60!black},
  breaklines=true,
  frame=single,
  columns=fullflexible,
  showstringspaces=false
}

% ─── Forest dirtree style ───────────────────────────────────────────────────
\forestset{
  dir tree/.style={
    for tree={
      font=\ttfamily\tiny,
      grow'=0,
      child anchor=west,
      parent anchor=south,
      anchor=west,
      calign=first,
      edge path={
        \noexpand\path [draw, -{Stealth[length=1.5pt]}]
        (!u.south west) +(4pt,0) |- node[fill,inner sep=0.6pt] {} (.child anchor)\forestoption{edge label};
      },
      before typesetting nodes={
        if n=1 {insert before={[,phantom]}} {}
      },
      fit=band,
      before computing xy={l=6pt},
      inner ysep=0pt,
    }
  }
}

% ─── Métadonnées ─────────────────────────────────────────────────────────────
\title[Bloc 2 — R-Type]{Bloc 2 — Concevoir une architecture logicielle}
\subtitle{Projet R-Type — Jeu multijoueur en C++}
\author{Dytoma BATOGOUMA, Baptiste CAMERLYNCK, Jonathan UFARTE,\newline Noe ROBERTIES, Marc GRIOCHE}
\date{16 février 2026}

% ─── Navigation simplifiée ───────────────────────────────────────────────────
% Keep only previous/next navigation arrows (remove all other navigation symbols)
\setbeamertemplate{navigation symbols}{
  \insertframenavigationsymbol
}

% ─── Footline personnalisée ──────────────────────────────────────────────────
% Remove authors from footer, keep short title, date, and slide numbers
\setbeamertemplate{footline}{
  \leavevmode%
  \hbox{%
    \begin{beamercolorbox}[wd=.5\paperwidth,ht=2.25ex,dp=1ex,left]{author in head/foot}%
      \usebeamerfont{author in head/foot}\hspace*{1em}\insertshorttitle
    \end{beamercolorbox}%
    \begin{beamercolorbox}[wd=.4\paperwidth,ht=2.25ex,dp=1ex,center]{title in head/foot}%
      \usebeamerfont{title in head/foot}\insertdate
    \end{beamercolorbox}%
    \begin{beamercolorbox}[wd=.1\paperwidth,ht=2.25ex,dp=1ex,right]{date in head/foot}%
      \usebeamerfont{date in head/foot}\insertframenumber{} / \inserttotalframenumber\hspace*{1em}
    \end{beamercolorbox}%
  }%
  \vskip0pt%
}

% ══════════════════════════════════════════════════════════════════════════════
\begin{document}

% ─── PAGE DE TITRE ───────────────────────────────────────────────────────────
\begin{frame}
  \titlepage
\end{frame}

% ─── SOMMAIRE ────────────────────────────────────────────────────────────────
\begin{frame}{Sommaire}
  \tableofcontents
\end{frame}

% ══════════════════════════════════════════════════════════════════════════════
%  A3 — VEILLE TECHNOLOGIQUE
% ══════════════════════════════════════════════════════════════════════════════
\section{A3 — Veille technologique}

% ──────────────────────────────────────────────────────────────────────────────
% SLIDE 1 : B2-C06.1 — Étude comparative
% ──────────────────────────────────────────────────────────────────────────────
\begin{frame}{B2-C06.1 — Étude comparative}
\footnotesize
Le projet R-Type a fait l'objet d'une étude comparative qualitative des technologies disponibles.

\vspace{0.2em}
\begin{adjustbox}{max width=\textwidth}
\begin{tabular}{l c c c c}
\toprule
\textbf{Domaine} & \textbf{Option 1} & \textbf{Option 2} & \textbf{Option 3} & \textbf{Choix retenu} \\
\midrule
Langage       & C++          & Rust        & C\#          & \textbf{C++} \\ \hline
Graphisme     & SDL2         & SFML        & Raylib       & \textbf{SDL2} \\ \hline
Réseau        & TCP          & UDP brut    & UDP + fiabilité & \textbf{UDP + fiabilité} \\ \hline
Sérialisation & JSON         & Protobuf    & Binaire custom  & \textbf{Binaire custom} \\ \hline
Architecture  & Monolithe    & Micro-services & ECS + OOP  & \textbf{ECS + OOP} \\ \hline
Build system  & Makefile     & Meson       & CMake        & \textbf{CMake} \\
\bottomrule
\end{tabular}
\end{adjustbox}

\vspace{0.3em}
Chaque choix évalué selon : \textbf{performance}, \textbf{latence}, \textbf{complexité}, \textbf{scalabilité} et \textbf{sécurité}.

\vspace{0.2em}
\end{frame}

% ──────────────────────────────────────────────────────────────────────────────
% SLIDE 2 : B2-C06.2 — Identification des technologies
% ──────────────────────────────────────────────────────────────────────────────
\begin{frame}{B2-C06.2 — Identification des technologies}
\footnotesize
Technologies identifiées et retenues pour le projet :

\vspace{0.2em}
\begin{adjustbox}{max width=\textwidth}
\begin{tabular}{l l l}
\toprule
\textbf{Catégorie} & \textbf{Technologie} & \textbf{Rôle} \\
\midrule
Langage        & C++ (C++17)           & Logique de jeu, systèmes bas niveau \\ \hline
Graphisme      & SDL2                  & Rendu 2D, gestion fenêtre, entrées \\ \hline
Réseau         & BSD Sockets (UDP)     & Communication temps réel client/serveur \\ \hline
Base de données & SQLite3              & Authentification, persistance XP joueurs \\ \hline
Sécurité       & libsodium             & Hachage de mots de passe (Argon2) \\ \hline
Configuration  & nlohmann/json         & Chargement des niveaux (JSON) \\ \hline
Build          & CMake                 & Compilation cross-platform \\ \hline
Conteneurisation & Docker \& Docker Compose & Déploiement serveur/client \\ \hline
Monitoring     & Prometheus + Grafana  & Métriques serveur en temps réel \\ \hline
Débogage       & Valgrind, GDB         & Détection fuites mémoire, debug runtime \\
\bottomrule
\end{tabular}
\end{adjustbox}

\vspace{0.3em}
Stack technologique offrant contrôle complet sur les performances, maintenabilité et sécurité.
\end{frame}

% ──────────────────────────────────────────────────────────────────────────────
% SLIDE 3 : B2-C07.1 — Étude sécurité informatique
% ──────────────────────────────────────────────────────────────────────────────
\begin{frame}{B2-C07.1 — Étude sécurité informatique}
\footnotesize
Étude des failles de sécurité des technologies utilisées :

\vspace{0.2em}
\begin{adjustbox}{max width=\textwidth}
\begin{tabular}{l l l}
\toprule
\textbf{Menace identifiée} & \textbf{Technologie} & \textbf{Mitigation implémentée} \\
\midrule
Buffer overflow UDP         & BSD Sockets           & Validation taille paquets, bounds checking \\ \hline
Injection SQL               & SQLite3               & Requêtes préparées (\texttt{sqlite3\_prepare\_v2}) \\ \hline
Mots de passe en clair      & Authentification      & Hachage Argon2id via \textbf{libsodium} \\ \hline
Flooding / DDoS             & Serveur UDP           & Rate limiting, quotas de connexion \\ \hline
Triche (speed hack)         & Gameplay              & Serveur autoritaire, validation serveur \\ \hline
Paquets falsifiés           & Protocole binaire     & Checksums, numéros de séquence \\ \hline
État client manipulé        & Client                & Aucun état client n'est ``de confiance'' \\
\bottomrule
\end{tabular}
\end{adjustbox}

\vspace{0.3em}
\textbf{Principe clé :} Le client n'envoie que des \textit{inputs}, jamais de résultats. Le serveur est la seule source de vérité.
\end{frame}

% ──────────────────────────────────────────────────────────────────────────────
% SLIDE 4 : B2-C07.2 — Veille sécurité informatique
% ──────────────────────────────────────────────────────────────────────────────
\begin{frame}{B2-C07.2 — Veille sécurité informatique}
\footnotesize
Connaissance de l'actualité sécurité informatique appliquée au projet :

\begin{itemize}\setlength\itemsep{2pt}
  \item \textbf{SDL2 :} Suivi des CVE (ex. CVE-2022-4743 — heap overflow SDL2 < 2.26.0). Version maintenue à jour.
  \item \textbf{SQLite3 :} Vulnérabilités d'injection contournées par \textit{prepared statements} systématiques.
  \item \textbf{libsodium :} Argon2id recommandé par l'OWASP. Gestion auto du \textit{rehashing} (\texttt{needsRehash()}).
  \item \textbf{UDP / Réseau :} Risques de \textit{spoofing}, \textit{replay attacks}. Mitigation par numéros de séquence et validation stricte.
  \item \textbf{OWASP Top 10 :} A01:Broken Access Control, A03:Injection, A07:Security Misconfiguration.
  \item \textbf{C++ moderne :} \texttt{std::unique\_ptr}, \texttt{std::shared\_ptr} contre fuites mémoire et \textit{use-after-free}.
\end{itemize}

\vspace{0.2em}
\textbf{Pratique :} Veille régulière via NVD et bulletins de sécurité des bibliothèques utilisées.
\end{frame}

% ══════════════════════════════════════════════════════════════════════════════
%  A4 — PROTOTYPAGE
% ══════════════════════════════════════════════════════════════════════════════
\section{A4 — Prototypage}

% ──────────────────────────────────────────────────────────────────────────────
% SLIDE 5 : B2-C08.1 — Prototypes
% ──────────────────────────────────────────────────────────────────────────────
\begin{frame}{B2-C08.1 — Prototypes}
\footnotesize
Différents prototypes réalisés au cours du développement :

\begin{enumerate}\setlength\itemsep{2pt}
  \item \textbf{Prototype réseau TCP} : Communication client/serveur en TCP. Fonctionnel mais latence trop élevée.
  \item \textbf{Prototype réseau UDP brut} : Migration vers UDP pur. Gain de latence majeur, perte de fiabilité.
  \item \textbf{Prototype UDP + couche de fiabilité} : Système ACK/NACK sélectif sur UDP. Compromis optimal retenu.
  \item \textbf{Prototype ECS} : Implémentation initiale \textit{Array of Components}, puis migration vers \textit{Sparse Set}.
  \item \textbf{Prototype rendu SDL2} : Pipeline de rendu 2D avec parallax, sprites animés et gestion d'entrées.
  \item \textbf{Prototype système de lobby} : Architecture multi-lobby avec threads dédiés et file d'attente d'inputs.
\end{enumerate}

\vspace{0.2em}
Chaque prototype a été évalué avant d'être intégré ou remplacé dans le projet final.
\end{frame}

% ──────────────────────────────────────────────────────────────────────────────
% SLIDE 6 : B2-C08.2 — Prototypes - comparatif
% ──────────────────────────────────────────────────────────────────────────────
\begin{frame}{B2-C08.2 — Prototypes — comparatif}
\footnotesize
Avantages et inconvénients de chaque prototype :

\vspace{0.2em}
\begin{adjustbox}{max width=\textwidth}
\begin{tabular}{l p{5.0cm} p{5.0cm}}
\toprule
\textbf{Prototype} & \textbf{Avantages} & \textbf{Inconvénients} \\
\midrule
TCP                  & Fiable, ordonné, simple      & Latence élevée (head-of-line blocking) \\ \hline
UDP brut             & Très faible latence           & Perte de paquets sans récupération \\ \hline
UDP + fiabilité      & Latence faible + fiabilité sélective & Complexité modérée à implémenter \\ \hline
ECS Array of Components & Cache-friendly, itération rapide & Rigide, mapping d'index complexe \\ \hline
ECS Sparse Set       & Flexible, ajout/suppression dynamique & Légèrement plus complexe \\ \hline
SDL2                 & Cross-platform, léger, accès bas niveau & API C-style, moins abstrait \\ \hline
SFML                 & API C++ OOP intuitive         & Moins de contrôle bas niveau \\
\bottomrule
\end{tabular}
\end{adjustbox}

\vspace{0.3em}
\textbf{Choix finaux :} UDP + couche de fiabilité, Sparse Set ECS, SDL2 — meilleur compromis performance/flexibilité/contrôle.
\end{frame}

% ──────────────────────────────────────────────────────────────────────────────
% SLIDE 7 : B2-C09.1 — Architecture
% ──────────────────────────────────────────────────────────────────────────────
\begin{frame}[fragile]{B2-C09.1 — Architecture}
\footnotesize
\textbf{Architecture ECS + Threads \& Réseau UDP :}

\vspace{0.1em}
\begin{columns}[T]
\begin{column}{0.32\textwidth}
\textbf{Patron ECS :}
\begin{itemize}\setlength\itemsep{1pt}
  \item \textbf{Entity} : identifiant unique
  \item \textbf{Component} : données pures (Position, Velocity, Health\ldots)
  \item \textbf{System} : logique (Movement, Collision, Weapon\ldots)
\end{itemize}

\vspace{0.3em}
\textbf{Threads :}
\begin{itemize}\setlength\itemsep{1pt}
  \item \textbf{Serveur} : Main (I/O) + 1 Thread/Lobby.
  \item \textbf{Client} : Main (Rendu) + Network.
\end{itemize}
\end{column}

\begin{column}{0.66\textwidth}
\centering
\begin{tikzpicture}[scale=0.92, transform shape]
  \tikzset{
    node distance=0.3cm,
    font=\tiny,
    box/.style={draw, rounded corners, align=center, fill=white},
    thread/.style={draw, dashed, rounded corners, fill=rtypelight, inner sep=2pt, align=center},
    packet/.style={font=\tiny\itshape, text=rtypegray, align=center}
  }

  % Server Container
  \node[box, fill=rtypeorange!20, minimum width=3.2cm, minimum height=3.2cm, label={[text=rtypeorange]north:\textbf{Serveur}}] (server) {};

  % Server Threads
  \node[thread, below=0.35cm of server.north, minimum width=2.8cm] (srv_main) {\textbf{Main Thread}\\(I/O, Auth)};
  
  \node[thread, below=0.2cm of srv_main, minimum width=1.3cm, xshift=-0.7cm] (lobby1) {\textbf{Lobby 1}\\(Thread)};
  \node[thread, below=0.2cm of srv_main, minimum width=1.3cm, xshift=0.7cm] (lobby2) {\textbf{Lobby 2}\\(Thread)};

  % Client Container
  \node[box, fill=rtypeblue!20, minimum width=3.2cm, minimum height=2.4cm, label={[text=rtypeblue]north:\textbf{Client}}, left=1.8cm of server] (client) {};

  % Client Threads
  \node[thread, below=0.35cm of client.north, minimum width=2.8cm] (cli_main) {\textbf{Main Thread}\\(Display, Input)};
  \node[thread, below=0.2cm of cli_main, minimum width=2.8cm] (cli_net) {\textbf{Network Thread}\\(Select/Poll)};

  % UDP Arrows
  \draw[-{Stealth}, thick, rtypeblue] ([yshift=0.4cm]client.east) -- node[above, packet] {INPUT, CONNECT,\\PING} ([yshift=0.4cm]server.west);
  
  \draw[-{Stealth}, thick, rtypeorange] ([yshift=-0.4cm]server.west) -- node[below, packet] {GAME\_STATE,\\SPAWN\_ENTITY} ([yshift=-0.4cm]client.east);

  % Database - Moved below server to save horizontal space or tightly right
  \node[box, fill=rtypegreen!20, right=0.2cm of server, anchor=west, minimum width=1.2cm] (db) {SQLite3};
  \draw[{Stealth}-{Stealth}, thick] (srv_main.east) -| (db.north);

\end{tikzpicture}
\end{column}
\end{columns}
\end{frame}

% ──────────────────────────────────────────────────────────────────────────────
% SLIDE 8 : B2-C09.2 — Intégration technique
% ──────────────────────────────────────────────────────────────────────────────
\begin{frame}{B2-C09.2 — Intégration technique}
\footnotesize
Justification de l'architecture et intégration dans l'écosystème technique :

\begin{itemize}\setlength\itemsep{2pt}
  \item \textbf{CMake} : Build cross-platform (Linux, Windows) — un \texttt{CMakeLists.txt} racine orchestrant \texttt{client/}, \texttt{server/}, \texttt{shared/}.
  \item \textbf{Docker} : Conteneurisation via \texttt{Dockerfile.server} et \texttt{Dockerfile.client}, orchestration \texttt{docker-compose.yml}.
  \item \textbf{Monitoring} : Stack Prometheus + Grafana pour métriques serveur (latence, connexions, ticks/s).
  \item \textbf{Shared library} : Code partagé (ECS, réseau, sérialisation) compilé en bibliothèque commune.
  \item \textbf{Modularité} : ECS permet d'ajouter composants/systèmes sans modifier l'existant (Open/Closed).
  \item \textbf{Scalabilité} : Système de lobbies avec threads dédiés et pool de workers.
\end{itemize}

\vspace{0.2em}
Écosystème \textbf{reproductible} (Docker), \textbf{observable} (Grafana) et \textbf{extensible} (ECS).
\end{frame}

% ══════════════════════════════════════════════════════════════════════════════
%  A5 — RÉSOLUTION DE PROBLÈMES ALGORITHMIQUES COMPLEXES
% ══════════════════════════════════════════════════════════════════════════════
\section{A5 — Résolution de problèmes algorithmiques complexes}

% ──────────────────────────────────────────────────────────────────────────────
% SLIDE 9 : B2-C10.1 — Implémentation de l'architecture
% ──────────────────────────────────────────────────────────────────────────────
\begin{frame}[fragile]{B2-C10.1 — Implémentation de l'architecture}
\footnotesize
L'implémentation respecte l'architecture ECS + Client/Serveur :

\begin{columns}[T]
\begin{column}{0.48\textwidth}
\textbf{Côté Shared (ECS) :}
\begin{lstlisting}
// Registry : coeur de l'ECS
class Registry {
  Entity create_entity();
  void kill_entity(Entity e);
  template<typename C>
  void emplace(Entity e, Args&&... args);
  template<typename... Cs>
  View<Cs...> view();
};
\end{lstlisting}

\textbf{ComponentStorage} : Sparse Set avec \texttt{dense\_entities}, \texttt{dense\_data}, \texttt{sparse} — accès $O(1)$.
\end{column}

\begin{column}{0.48\textwidth}
\textbf{Côté Serveur :}
\begin{lstlisting}
class RTypeServer : public RTypeNetwork {
  void handleConnect(...);
  void handleInput(...);
  void handleAuthRequest(...);
  LobbyManager _lobbyManager;
  std::unique_ptr<SqlDB::DB> _db;
};
\end{lstlisting}

\textbf{GameInstance} orchestre : Core, Player, Entities, Physics, Serialization, Conditions.
\end{column}
\end{columns}
\end{frame}

% ──────────────────────────────────────────────────────────────────────────────
% SLIDE 10 : B2-C10.2 — Bonnes pratiques de développement
% ──────────────────────────────────────────────────────────────────────────────
\begin{frame}[fragile]{B2-C10.2 — Bonnes pratiques de développement}
\footnotesize

\begin{columns}[T]
\begin{column}{0.48\textwidth}
\textbf{RAII et Smart Pointers :}
\begin{lstlisting}
// Ownership clair
std::unique_ptr<GameInstance> gameInstance;
std::unique_ptr<SqlDB::DB> _db;
std::unique_ptr<UdpSocket> _socket;
\end{lstlisting}

\textbf{Design Patterns :}
\begin{itemize}\setlength\itemsep{1pt}
  \item \textbf{ECS} : séparation données/logique
  \item \textbf{Factory} : \texttt{EntityFactory}
  \item \textbf{Observer} : dispatch handlers réseau
  \item \textbf{Template Method} : \texttt{RTypeNetwork}
\end{itemize}
\end{column}

\begin{column}{0.48\textwidth}
\textbf{Paradigmes utilisés :}
\begin{itemize}\setlength\itemsep{1pt}
  \item \textbf{OOP} : héritage, polymorphisme
  \item \textbf{Générique} : templates C++ (\texttt{ComponentStorage<T>}, \texttt{View<Cs...>})
  \item \textbf{Data-Oriented} : ECS cache-friendly
\end{itemize}

\vspace{0.2em}
\textbf{Qualité du code :}
\begin{itemize}\setlength\itemsep{1pt}
  \item Documentation Doxygen
  \item Pre-commit hooks
  \item Valgrind pour les fuites mémoire
\end{itemize}
\end{column}
\end{columns}
\end{frame}

% ──────────────────────────────────────────────────────────────────────────────
% SLIDE 11 : B2-C11.1 — Organisation du code
% ──────────────────────────────────────────────────────────────────────────────
\begin{frame}[fragile]{B2-C11.1 — Organisation du code}
\scriptsize

\begin{columns}[T]
\begin{column}{0.42\textwidth}
\textbf{Arborescence du projet :}

\vspace{0.1em}
\begin{forest}
  dir tree
  [\textbf{r-type/}
    [\textcolor{rtypeblue}{\textbf{client/src/}}
      [game\_handler/]
      [managers/]
      [network/]
      [render/]
      [ui/]
      [prediction/]
    ]
    [\textcolor{rtypeorange}{\textbf{server/}}
      [src/ {\tiny Handle*.cpp}]
      [db/ {\tiny DB / PasswordManager}]
      [metrics/]
    ]
    [\textcolor{rtypegreen}{\textbf{shared/src/}}
      [ecs/ {\tiny Registry / Components / Systems}]
      [game\_instance/ {\tiny 7 sous-modules}]
      [entities/]
      [levels/]
    ]
    [Docker/]
    [monitoring/]
  ]
\end{forest}
\end{column}

\begin{column}{0.55\textwidth}
\textbf{Gestion de version — Git :}
\begin{itemize}\setlength\itemsep{1pt}
  \item \textbf{73 branches} feature créées au total
  \item Nommage par issue : \texttt{XX-description}\\
        {\scriptsize Ex: \texttt{3-basic-ecs-structure}, \texttt{29-player-synchronisation-rollback}}
  \item Branche \texttt{main} protégée — merge après revue
  \item Pre-commit hooks (\texttt{.pre-commit-config.yaml}) pour la qualité automatique
\end{itemize}

\vspace{0.2em}
\textbf{Persistance des données :}

\vspace{0.1em}
\begin{adjustbox}{max width=\columnwidth}
\begin{tabular}{l l}
\toprule
\textbf{Donnée} & \textbf{Solution} \\
\midrule
Auth. joueurs & SQLite3 (requêtes préparées) \\
XP joueurs    & SQLite3 \\
Config. niveaux & JSON (nlohmann/json) \\
État réseau   & Binaire custom (UDP) \\
État en mémoire & \texttt{std::vector}, \texttt{unordered\_map} \\
\bottomrule
\end{tabular}
\end{adjustbox}

\end{column}
\end{columns}
\end{frame}

% ──────────────────────────────────────────────────────────────────────────────
% SLIDE 12 : B2-C11.2 — Segmentation du code
% ──────────────────────────────────────────────────────────────────────────────
\begin{frame}[fragile]{B2-C11.2 — Segmentation du code}
\scriptsize

\begin{columns}[T]
\begin{column}{0.44\textwidth}
\textbf{GameInstance — 7 sous-modules :}

\vspace{0.1em}
\begin{forest}
  dir tree
  [\textbf{game\_instance/}
    [\textcolor{rtypeblue}{Core.cpp} {\tiny état / cycle de vie}]
    [\textcolor{rtypeblue}{Player.cpp} {\tiny joueurs / inputs}]
    [\textcolor{rtypeblue}{Entities.cpp} {\tiny spawn/destroy}]
    [\textcolor{rtypeblue}{Physics.cpp} {\tiny physique / collisions}]
    [\textcolor{rtypeblue}{Serialization.cpp} {\tiny sérialisation}]
    [\textcolor{rtypeblue}{Conditions.cpp} {\tiny win/lose}]
    [\textcolor{rtypeblue}{InputProcessor.cpp} {\tiny inputs}]
  ]
\end{forest}

\vspace{0.2em}
\textbf{Serveur — 1 fichier par message :}

\vspace{0.1em}
\begin{forest}
  dir tree
  [\textbf{server/src/}
    [\textcolor{rtypeorange}{HandleConnect.cpp}]
    [\textcolor{rtypeorange}{HandleInput.cpp}]
    [\textcolor{rtypeorange}{HandleAuth.cpp}]
    [\textcolor{rtypeorange}{HandleCreateLobby.cpp}]
    [\textcolor{rtypeorange}{\ldots{} (11 handlers total)}]
  ]
\end{forest}
\end{column}

\begin{column}{0.53\textwidth}
\textbf{Segmentation ECS :}

\vspace{0.1em}
\begin{forest}
  dir tree
  [\textbf{shared/src/ecs/}
    [Registry.hpp {\tiny cœur ECS}]
    [ComponentStorage.hpp {\tiny Sparse Set}]
    [\textcolor{rtypegreen}{components/} {\tiny 30+ structs}
      [{Position{,} Velocity{,} Health\ldots}]
    ]
    [\textcolor{rtypegreen}{systems/} {\tiny 15+ systèmes}
      [{Movement{,} Collision{,} Weapon\ldots}]
    ]
  ]
\end{forest}

\vspace{0.2em}
\textbf{Avantages :}
\begin{itemize}\setlength\itemsep{0pt}
  \item \textbf{Single Responsibility} : 1 fichier = 1 responsabilité
  \item \textbf{Isolation} : modifier la physique ne touche pas la sérialisation
  \item \textbf{Testable} : chaque module testable en isolation
  \item \textbf{Scalable} : ajouter un composant = ajouter un fichier
\end{itemize}
\end{column}
\end{columns}
\end{frame}

% ──────────────────────────────────────────────────────────────────────────────
% SLIDE 13 : B2-C12.1 — Implémentation algorithmique
% ──────────────────────────────────────────────────────────────────────────────
\begin{frame}{B2-C12.1 — Implémentation algorithmique}
\footnotesize
Algorithmes existants et optimaux implémentés :

\begin{itemize}\setlength\itemsep{2pt}
  \item \textbf{Spatial Hashing} pour la détection de collisions :
    \begin{itemize}\setlength\itemsep{0pt}
      \item Coût moyen $O(1)$ par requête, adapté aux entités mobiles en 2D.
      \item Préféré au Quadtree (plus adapté aux environnements statiques).
    \end{itemize}
  \item \textbf{Client-Side Prediction} :
    \begin{itemize}\setlength\itemsep{0pt}
      \item Simulation locale immédiate pour masquer la latence.
      \item Standard dans les jeux multijoueurs (Quake, Source Engine).
    \end{itemize}
  \item \textbf{Server Reconciliation} :
    \begin{itemize}\setlength\itemsep{0pt}
      \item Replay des inputs non confirmés après réception de l'état serveur.
      \item Interpolation linéaire pour lisser les corrections.
    \end{itemize}
  \item \textbf{Sparse Set} pour le stockage ECS :
    \begin{itemize}\setlength\itemsep{0pt}
      \item Accès, insertion et suppression en $O(1)$.
      \item Itération cache-friendly via vecteurs denses contigus.
    \end{itemize}
\end{itemize}
\end{frame}

% ──────────────────────────────────────────────────────────────────────────────
% SLIDE 14 : B2-C12.2 — Implémentation d'algorithmes originaux
% ──────────────────────────────────────────────────────────────────────────────
\begin{frame}{B2-C12.2 — Implémentation d'algorithmes originaux}
\footnotesize

\begin{columns}[T]
\begin{column}{0.54\textwidth}
\textbf{Protocole binaire custom :}

\vspace{0.2em}
\centering
\begin{tikzpicture}[
  field/.style={draw, minimum height=0.5cm, font=\tiny\ttfamily, fill=#1, text=white, align=center},
  label/.style={font=\tiny, text=rtypegray, align=center}
]
  \node[field=rtypeblue, minimum width=0.85cm] (t) at (0,0) {Type};
  \node[field=rtypegreen, minimum width=0.85cm, right=0pt of t] (v) {Ver};
  \node[field=rtypeorange, minimum width=1.1cm, right=0pt of v] (s) {Seq};
  \node[field=rtypeblue!70, minimum width=1.4cm, right=0pt of s] (p) {PlayerID};
  \node[field=rtypegray, minimum width=0.9cm, right=0pt of p] (l) {Len};
  \node[field=rtypedark, minimum width=1.6cm, right=0pt of l] (d) {Payload};

  \node[label, below=1pt of t] {1B};
  \node[label, below=1pt of v] {1B};
  \node[label, below=1pt of s] {2B};
  \node[label, below=1pt of p] {4B};
  \node[label, below=1pt of l] {2B};
  \node[label, below=1pt of d] {NB};
\end{tikzpicture}

\vspace{0.3em}
\raggedright
\begin{itemize}\setlength\itemsep{1pt}
  \item Sérialisation avec gestion du \textit{byte swapping}
  \item \textbf{ACK/NACK sélectif} : seuls les messages critiques sont acquittés
  \item \textbf{Agrégation} : petites mises à jour regroupées
\end{itemize}
\end{column}

\begin{column}{0.43\textwidth}
\textbf{Autres solutions originales :}

\begin{itemize}\setlength\itemsep{1pt}
  \item \textbf{Niveaux par JSON :}
    \begin{itemize}\setlength\itemsep{0pt}
      \item Chargement dynamique (gravité, ennemis, boss, plateformes)
      \item Nouveaux niveaux sans recompilation
    \end{itemize}
  \item \textbf{ECS View optimisé :}
    \begin{itemize}\setlength\itemsep{0pt}
      \item Itération multi-composants via intersection des sparse sets
      \item Sélection auto du plus petit set
    \end{itemize}
  \item \textbf{Messages temps réel :} envoyés sans garantie (positions, inputs) pour minimiser la latence
\end{itemize}
\end{column}
\end{columns}
\end{frame}

% ══════════════════════════════════════════════════════════════════════════════
%  A6 — DÉFINITION DES MODÈLES DE DONNÉES
% ══════════════════════════════════════════════════════════════════════════════
\section{A6 — Définition des modèles de données}

% ──────────────────────────────────────────────────────────────────────────────
% SLIDE 15 : B2-C13.1 — Persistance des données
% ──────────────────────────────────────────────────────────────────────────────
\begin{frame}{B2-C13.1 — Persistance des données}
\footnotesize
Les choix de persistance sont cohérents avec les besoins du projet :

\vspace{0.2em}
\begin{adjustbox}{max width=\textwidth}
\begin{tabular}{l l l l}
\toprule
\textbf{Donnée} & \textbf{Solution} & \textbf{Contrainte} & \textbf{Justification} \\
\midrule
Authentification  & SQLite3         & Sécurité, intégrité      & Requêtes préparées, ACID \\
XP joueurs        & SQLite3         & Persistance              & Transactions atomiques \\
Niveaux de jeu    & Fichiers JSON   & Lisibilité               & Édition sans recompilation \\
État de jeu       & Mémoire (RAM)   & Performance temps réel   & Accès $< 1\mu s$ \\
Communication     & Binaire (UDP)   & Latence minimale         & Compact, pas de parsing lourd \\
\bottomrule
\end{tabular}
\end{adjustbox}

\vspace{0.3em}
\textbf{Sécurité :}
\begin{itemize}\setlength\itemsep{1pt}
  \item Mots de passe hachés avec \textbf{Argon2id} (libsodium) — jamais stockés en clair.
  \item \texttt{PasswordManager::needsRehash()} pour adapter les paramètres de sécurité.
  \item Backup de la base via \texttt{DB::backup\_to\_file()}.
\end{itemize}
\end{frame}

% ──────────────────────────────────────────────────────────────────────────────
% SLIDE 16 : B2-C13.2 — Persistance des données - comparatif
% ──────────────────────────────────────────────────────────────────────────────
\begin{frame}{B2-C13.2 — Persistance des données — comparatif}
\footnotesize
Comparaison des solutions de persistance envisagées :

\vspace{0.2em}
\textbf{Base de données :}
\begin{adjustbox}{max width=\textwidth}
\begin{tabular}{l c c c c}
\toprule
 & \textbf{SQLite3} & \textbf{PostgreSQL} & \textbf{MongoDB} & \textbf{Fichiers texte} \\
\midrule
Déploiement     & Embarqué, simple  & Serveur dédié   & Serveur dédié  & Trivial \\
Performance     & Bonne (local)     & Excellente      & Bonne          & Variable \\
Requêtes SQL    & Oui               & Oui             & Non (NoSQL)    & Non \\
Dépendances     & Aucune (intégré)  & Lourde          & Lourde         & Aucune \\
\textbf{Choix}  & \textbf{$\checkmark$ Retenu} & Non retenu & Non retenu & Non retenu \\
\bottomrule
\end{tabular}
\end{adjustbox}

\vspace{0.2em}
\textbf{Format de configuration :}
\begin{adjustbox}{max width=\textwidth}
\begin{tabular}{l c c c c}
\toprule
 & \textbf{JSON} & \textbf{YAML} & \textbf{TOML} & \textbf{Binaire} \\
\midrule
Lisibilité    & Bonne     & Excellente & Bonne    & Nulle \\
Parsing C++   & nlohmann  & yaml-cpp   & toml++   & Custom \\
Écosystème    & Très large & Large     & Modéré   & Aucun \\
\textbf{Choix} & \textbf{$\checkmark$ Retenu} & Non retenu & Non retenu & Non retenu \\
\bottomrule
\end{tabular}
\end{adjustbox}

\vspace{0.2em}
\textbf{Justification :} SQLite3 embarqué (zéro config), JSON supporté nativement par nlohmann/json.
\end{frame}

% ──────────────────────────────────────────────────────────────────────────────
% SLIDE 17 : B2-C14.1 — Structures de données
% ──────────────────────────────────────────────────────────────────────────────
\begin{frame}{B2-C14.1 — Structures de données}
\footnotesize

\begin{columns}[T]
\begin{column}{0.50\textwidth}
\textbf{Structures utilisées :}

\vspace{0.1em}
\begin{adjustbox}{max width=\columnwidth}
\begin{tabular}{l l l}
\toprule
\textbf{Structure} & \textbf{Usage} & \textbf{Complexité} \\
\midrule
\texttt{std::vector}         & Composants ECS   & $O(1)$ accès \\
\texttt{std::unordered\_map} & Entité $\to$ idx  & $O(1)$ lookup \\
\texttt{std::queue}          & Messages réseau   & $O(1)$ push/pop \\
\texttt{std::unique\_ptr}    & Ownership RAII    & Zéro overhead \\
Sparse Set                    & ComponentStorage  & $O(1)$ CRUD \\
\texttt{std::atomic}         & Flags threads     & Lock-free \\
\bottomrule
\end{tabular}
\end{adjustbox}
\end{column}

\begin{column}{0.47\textwidth}
\textbf{Sparse Set — visualisation :}

\vspace{0.1em}
\centering
\begin{tikzpicture}[
  cell/.style={draw, minimum width=0.55cm, minimum height=0.4cm, font=\tiny\ttfamily, align=center},
  lbl/.style={font=\tiny\bfseries, text=rtypeblue},
  arr/.style={-{Stealth[length=1.5pt]}, rtypegray, thin},
]
  \node[lbl] at (-0.7, 1.8) {sparse};
  \node[cell, fill=rtypelight] (s0) at (0, 1.8) {1};
  \node[cell, fill=rtypelight] (s1) at (0.6, 1.8) {\textendash};
  \node[cell, fill=rtypelight] (s2) at (1.2, 1.8) {0};
  \node[cell, fill=rtypelight] (s3) at (1.8, 1.8) {\textendash};
  \node[cell, fill=rtypelight] (s4) at (2.4, 1.8) {2};
  \node[font=\tiny, below=0pt of s0] {id:0};
  \node[font=\tiny, below=0pt of s2] {id:2};
  \node[font=\tiny, below=0pt of s4] {id:4};

  \node[lbl] at (-0.7, 0.8) {dense\_e};
  \node[cell, fill=rtypeblue!20] (de0) at (0, 0.8) {E2};
  \node[cell, fill=rtypeblue!20] (de1) at (0.6, 0.8) {E0};
  \node[cell, fill=rtypeblue!20] (de2) at (1.2, 0.8) {E4};

  \node[lbl] at (-0.7, 0.2) {dense\_d};
  \node[cell, fill=rtypegreen!20] (dd0) at (0, 0.2) {Pos};
  \node[cell, fill=rtypegreen!20] (dd1) at (0.6, 0.2) {Pos};
  \node[cell, fill=rtypegreen!20] (dd2) at (1.2, 0.2) {Pos};

  \draw[arr] (s0.south) -- (de1.north);
  \draw[arr] (s2.south) -- (de0.north);
  \draw[arr] (s4.south) -- (de2.north);
\end{tikzpicture}

\vspace{0.1em}
{\tiny\texttt{sparse[entity\_id]} $\to$ index dense.\\
Mémoire contiguë = cache-friendly.}
\end{column}
\end{columns}
\end{frame}

% ──────────────────────────────────────────────────────────────────────────────
% SLIDE 18 : B2-C14.2 — Structures de données - justifications
% ──────────────────────────────────────────────────────────────────────────────
\begin{frame}{B2-C14.2 — Structures de données — justifications}
\footnotesize

\begin{columns}[T]
\begin{column}{0.48\textwidth}
\textbf{Performance :}
\begin{itemize}\setlength\itemsep{1pt}
  \item \texttt{std::vector} : mémoire contiguë, cache CPU optimisé.
  \item Sparse Set : $O(1)$ pour \texttt{has()}, \texttt{get()}, \texttt{emplace()}, \texttt{remove()}. Zéro allocation à l'itération.
  \item \texttt{std::queue} : push/pop $O(1)$, pattern producteur/consommateur (I/O $\to$ workers).
\end{itemize}

\textbf{Maintenabilité :}
\begin{itemize}\setlength\itemsep{1pt}
  \item \texttt{unique\_ptr} : ownership clair, pas de \texttt{delete} manuel.
  \item Templates : \texttt{ComponentStorage<T>} générique pour tout composant.
\end{itemize}
\end{column}

\begin{column}{0.48\textwidth}
\textbf{Évolutivité :}
\begin{itemize}\setlength\itemsep{1pt}
  \item Nouveau composant ECS = créer un \texttt{struct}, l'enregistrer. Aucune modification du Registry.
  \item Nouveau système = une classe avec \texttt{update()}, itérer via \texttt{registry.view<Cs...>()}.
  \item Protocole réseau : jusqu'à 255 types de messages (\texttt{uint8\_t}).
\end{itemize}

\vspace{0.2em}
\textbf{Conclusion :}\\
Accès $O(1)$ pour les opérations critiques. Extension simple grâce aux templates et au RAII.
\end{column}
\end{columns}
\end{frame}

% ─── SLIDE FINALE ────────────────────────────────────────────────────────────
\begin{frame}
\centering
\vspace{2cm}
{\Huge\bfseries Merci pour votre attention}

\vspace{1cm}
{\large Projet R-Type — Bloc 2}

\vspace{0.5cm}
{\small Questions ?}
\end{frame}

\end{document}